\section{Particle Decay and Rockets}
\subsection{a}
\paragraph{Problem:} The relative velocity $v$ between two particles can be determined from the four-scalar product of their four-momenta,
\begin{equation*}
    P_1 \cdot P_2 = m_1 m_2 c^2 \gamma(v)
\end{equation*}
(where $\gamma(v) \equiv  1/\sqrt{1 - v^2 / c^2} $). Show this formula and motivate well.

\paragraph{Solution:} Let the first particle be in rest. Letting $u_1$ and $u_2$ be the four-velocities then
\begin{align*}
    u_1 &= \gamma(0) (c, \vec{0})\\
    u_2 &= \gamma(v) (c, \vec{v}).
\end{align*}
Since $\gamma(0) = 1$ we can write the four-momentum as
\begin{align*}
    P_1 &= m_1 u_1 = m_1 (c, \vec{0})\\
    P_2 &= m_2 u_2 = m_2 \gamma(v) (c, \vec{v}).
\end{align*}
The calculation is therefore
\begin{align*}
    P_1 \cdot P_2 &= m_1 m_2 \gamma(v) (c, \vec{0}) \cdot (c, \vec{v}) = m_1 m_2 \gamma(v) (c^2 - \vec{0} \cdot \vec{v}) = m_1 m_2 c^2 \gamma(v)
\end{align*}
\qedsymbol

\subsection{b}
\paragraph{Problem:} Consider a particle A at rest, decaying into two particles B and C. Use suitable invariants to show that the velocity of B is determined by
\begin{equation*}
    \gamma(v_B) = \frac{m^2_A + m^2_B - m^2_C}{2 m_A m_B}
\end{equation*}

\paragraph{Solution:} Before the collision we have the four momentum
\begin{equation*}
    P_A = (E_A/c, \vec{0})
\end{equation*}
where $E_A = m_A c^2$ and thus
\begin{equation*}
    P_A = m_A c (1, \vec{0})
\end{equation*}

\subsection{c}
\paragraph{Problem:} Imagine that we can fuel a rocket with an ideal fuel, all of whose mass is converted to energy, leaving only massless exhaust that is released straight backwards. A rocket, that initially is at rest, is packed with ideal fuel. The fuel makes up a fraction $\alpha$ of the
loaded rocket's total mass. Derive the maximum velocity of the empty rocket after burning all the
fuel, as a function of $\alpha$. (\textit{Answer}: $v = c \frac{1- (1-\alpha)^2}{1+ (1-\alpha)^2}$ . \textit{Hint}: Consider the process as a ``decay''.)


\paragraph{Solution:}
